\chapter{导数：曲线在某一点的速度}

\section{过山车上的一个问题}

游乐场里最刺激的设施是过山车。车爬上一个又高又缓的坡，突然俯冲下去，再被甩上一个几乎竖直的环。坐在车里的人有一种特别清楚的感受：有些地方“平”，有些地方“陡”，有些地方正在“拐弯”，而速度好像每一瞬间都不一样。

现在我们把过山车的轨道画在坐标纸上，让它成为一条曲线 $y=f(x)$：横坐标 $x$ 是水平距离，纵坐标 $y$ 是离地高度。过山车手描述体验时说的是“这里最陡”“这里开始下坡”“这里几乎是平的”。这些话在数学上能精确表达吗？

再看一个更朴素的场景。一辆车沿公路行驶，里程表告诉你：上午 10 点，已经开了 120 千米；上午 11 点，开了 180 千米。于是你说：“这一小时的平均速度是 60 千米每小时。”这谁都同意。但交警拿出的测速照片显示：10 点 23 分 15 秒，这辆车正以 92 千米每小时的速度通过一个限速 80 的路段。请注意这个区别：交警说的不是“某段时间的平均速度”，而是“那一瞬的速度”。一瞬间时间间隔是零，路程增量也是零，速度难道不是 $\frac{0}{0}$ 吗？

这两个场景——轨道在某点的陡峭程度，汽车在某瞬的速度——其实是同一个数学问题。这个问题困扰了数学家很多年，而解决它的工具，就是本章的主角：导数。

\section{先猜一猜：瞬时速度到底存不存在}

在继续之前，你可以先凭直觉回答几个问题：

\begin{itemize}
  \item 我们说“某一时刻的速度是 92 千米每小时”，这句话真的有含义吗，还是只是口头方便？
  \item 一条弯弯曲曲的线，在它的某一个“点”上，谈“陡峭程度”合理吗？陡峭难道不是一整段路的事？
  \item 如果瞬时速度存在，为什么 $\frac{0}{0}$ 没有把它毁掉？
\end{itemize}

很多人凭直觉会说：“某一瞬间当然有一个确定的速度，我骑车的时候能感觉到。”另一些人会警觉起来：“速度的定义是路程除以时间，时间间隔为零，这个除法就没意义了。”两边都对：直觉对了一半，警觉也对了一半。数学家干的事，就是把两边的直觉都翻译成一个不矛盾的精确定义。上一章讲的极限，正是为这个翻译准备的工具。

这场翻译在十七世纪真的发生过，而且是两次。牛顿在研究行星为什么那样运行时，需要处理“每一瞬间的速度”，他管变化率叫“流数”；莱布尼茨则从几何问题出发，研究曲线的切线，发明了沿用至今的 $\frac{\mathrm{d}y}{\mathrm{d}x}$ 记号。两人各自独立地抓住了同一个思想，后来却为“谁先发明”争了很多年，甚至惊动了欧洲各国学会站队。这场争论在今天看来有些可惜：两人问的问题看似不同——一个来自运动，一个来自切线——答案却是同一个概念。这恰好说明这个概念有多自然：只要认真追问“变化”本身，谁都会走到它面前。

\section{第一次尝试：用很短的时间间隔}

我们先试着解决汽车的问题。设汽车行驶的路程随时间变化的规律是 $s(t)=t^2$（单位是千米和小时，数字当然不真实，但道理一样）。问题是：$t=2$ 这一时刻，瞬时速度是多少？

第一种尝试：取一个很小的间隔，比如 $\Delta t=0.1$ 小时，算平均速度
\[
\frac{s(2.1)-s(2)}{0.1}=\frac{4.41-4}{0.1}=4.1\ \text{千米每小时}.
\]
再取 $\Delta t=0.01$：
\[
\frac{s(2.01)-s(2)}{0.01}=\frac{4.0401-4}{0.01}=4.01\ \text{千米每小时}.
\]
$\Delta t=0.001$ 时答案是 $4.001$；$\Delta t=0.0001$ 时答案是 $4.0001$。你看，间隔越短，答案越逼近 4。于是我们猜：$t=2$ 时的瞬时速度是 4 千米每小时。

把这张“逼近过程”列成表格，会更清楚。我们还可以从左侧逼近——取 $t=2$ 之前的一小段，看看两边是不是奔向同一个数：

\begin{center}
\begin{tabular}{ccccc}
\toprule
$\Delta t$ & $1$ & $0.1$ & $0.01$ & $0.001$ \\
\midrule
从右侧：$\frac{s(2+\Delta t)-s(2)}{\Delta t}$ & $5$ & $4.1$ & $4.01$ & $4.001$ \\
从左侧：$\frac{s(2)-s(2-\Delta t)}{\Delta t}$ & $3$ & $3.9$ & $3.99$ & $3.999$ \\
\bottomrule
\end{tabular}
\end{center}

右列数字从上方逼近 4，左列数字从下方逼近 4，两边在同一个地方会师。这个“两边会师”很关键：如果只从一侧逼近，另一侧却奔向别的数，那就说明这一瞬的速度并不存在——本章后面思考题里的 $y=|x|$ 在原点处正是这种情况。

但请注意，我们始终不敢把 $\Delta t$ 真正取成 0，因为那样分子分母都是 0。我们真正用到的想法是：$\Delta t$ 可以任意小，而平均速度可以任意接近 4——这正是上一章“极限”的语言。瞬时速度不是某个很短的间隔算出来的平均速度，而是间隔趋于零时平均速度的极限。

用符号写出来：当 $\Delta t$ 趋于 0 时，$\frac{s(t+\Delta t)-s(t)}{\Delta t}$ 的极限，就是时刻 $t$ 的瞬时速度。

\begin{fanli}
有个流传很广的“证明”是错的：既然瞬时速度是时间间隔为零时的速度，那就把 $\Delta t=0$ 代进去，得到 $\frac{0}{0}$，所以瞬时速度不存在。

错在哪里？极限不是“代入”。$\Delta t$ 趋于 0 指的是它可以任意小，但始终不是 0。像 $\frac{s(t+\Delta t)-s(t)}{\Delta t}$ 这样的表达式，在 $\Delta t\neq 0$ 时常常可以先化简：例如 $\frac{(2+\Delta t)^2-4}{\Delta t}=\frac{4\Delta t+(\Delta t)^2}{\Delta t}=4+\Delta t$，化简之后再让 $\Delta t$ 趋于 0，就得到干干净净的 4。整个过程中我们一次也没有真正除以零。
\end{fanli}

\section{把同一个想法搬到曲线上：割线与切线}

现在我们转到几何。看抛物线 $y=x^2$，想在点 $(1,1)$ 处描述它的“陡峭程度”。一条直线的陡峭程度用斜率刻画，可这是一条弯的线，哪来的斜率？

模仿汽车问题的做法：不取一个点，取两个点。在曲线上另取一点 $(1+h,(1+h)^2)$，连接这两点的直线叫做割线，它的斜率是
\[
\frac{(1+h)^2-1}{h}=\frac{2h+h^2}{h}=2+h.
\]
取 $h=1$，割线斜率是 3；$h=0.1$，是 2.1；$h=0.01$，是 2.01。让 $h$ 趋于 0，割线的斜率趋于 2。在几何上这意味着：当第二个点沿着曲线滑向 $(1,1)$ 时，割线不断转动，最后“卡”在一条确定的直线上——这条直线就是曲线在该点的切线，它的斜率就是 2。

\begin{center}
\begin{tikzpicture}[scale=1.6]
  \draw[->] (-1.2,0) -- (2.8,0) node[right] {$x$};
  \draw[->] (0,-0.5) -- (0,4.6) node[above] {$y$};
  \draw[thick,domain=-1.15:2.05,samples=60] plot (\x,{\x*\x});
  \draw[dashed] (1,1) -- (2,4);
  \draw[blue!60!black,thick,domain=0.05:2.2,samples=60] plot (\x,{2*\x-1});
  \fill (1,1) circle (0.05) node[below left] {$(1,1)$};
  \fill (2,4) circle (0.05) node[above right] {$(2,4)$};
  \node[blue!60!black] at (2.3,2.4) {切线：斜率 $2$};
  \node at (1.6,3.4) {割线};
\end{tikzpicture}
\end{center}

现在两件事合一了：路程函数的瞬时速度，等于路程图像上切线的斜率。这个合一值得停一秒——代数里的“速度”和几何里的“斜率”，本来是两个世界的东西，在这里成了同一个数。这个数叫导数。

\begin{shuyu}{切线}
直观解释：在切点附近与曲线“贴合得最好”的那条直线——它不是“只碰一个点”的线（抛物线的对称轴也只碰顶点一个点，却不是切线），而是曲线在该点前进方向的化身。
正式定义：过曲线上定点与邻近点的割线，在邻近点趋于定点时的极限位置。
一个例子：$y=x^2$ 在 $(1,1)$ 处的切线是 $y=2x-1$，它从曲线下方整个托住抛物线。
\end{shuyu}

\begin{dingyi}[导数]
设函数 $f$ 在点 $x$ 附近有定义。如果极限
\[
f'(x)=\lim_{\Delta x\to 0}\frac{f(x+\Delta x)-f(x)}{\Delta x}
\]
存在，就称 $f$ 在 $x$ 处可导，这个极限叫做 $f$ 在 $x$ 处的导数，记作 $f'(x)$（物理学家也常写 $\frac{\mathrm{d}y}{\mathrm{d}x}$）。
\end{dingyi}

\begin{shuyu}{导数}
直观解释：函数在某一点“每走一小步横坐标，纵坐标大约变化多少”，也就是曲线在该点切线的斜率。
正式定义：平均变化率在区间长度趋于零时的极限。
一个例子：$f(x)=x^2$ 在 $x=1$ 处导数为 2，意味着在 $(1,1)$ 附近，横坐标每增加一点点，纵坐标大约增加它的两倍那么多。
\end{shuyu}

\section{亲手算几个导数：$x^2$、$x^3$ 和一次函数}

定义好不好用，算了才知道。我们用定义把三个函数的导数从头发明一遍。

先看 $f(x)=x^2$。对任意的 $x$：
\[
\frac{f(x+h)-f(x)}{h}=\frac{(x+h)^2-x^2}{h}=\frac{2xh+h^2}{h}=2x+h.
\]
让 $h$ 趋于 0，得到 $f'(x)=2x$。这一个公式就把整条抛物线上每一点的切线斜率都装进去了：$x=1$ 处斜率 2（正是刚才割线实验的结果）；$x=0$ 处斜率 0，说明抛物线在最底部是水平的；$x=-3$ 处斜率 $-6$，负数表示曲线正在下降。

再看 $f(x)=x^3$。这里要用一个中学代数恒等式：
\[
(x+h)^3=x^3+3x^2h+3xh^2+h^3.
\]
于是
\[
\frac{(x+h)^3-x^3}{h}=\frac{3x^2h+3xh^2+h^3}{h}=3x^2+3xh+h^2.
\]
$h$ 趋于 0 时，后两项消失，得到 $f'(x)=3x^2$。注意一个有趣的事：$3x^2$ 永远非负，这说明 $x^3$ 这条曲线处处都在“往上走或走平”，最平缓的地方是原点，那里斜率为 0——但注意，它不是山峰也不是谷底，只是一个“歇脚处”。你还可以比较两条曲线“变陡”的节奏：$x^2$ 的斜率 $2x$ 随 $x$ 均匀增长，$x^3$ 的斜率 $3x^2$ 却越涨越快。到 $x=10$ 时，抛物线的切线斜率是 20，而三次曲线已经达到 300。导数不只告诉你某一点有多陡，还告诉你“陡峭本身”如何变化——这个思想在后面会自己长成一个新的研究对象。

最后看一次函数 $f(x)=kx+b$：
\[
\frac{k(x+h)+b-(kx+b)}{h}=\frac{kh}{h}=k,
\]
导数就是 $k$ 本身。这当然应该如此：直线处处一样陡，它的“切线”就是它自己。

\begin{li}
火箭点火后离地高度 $h(t)=5t^2$ 米，求 $t=10$ 秒时的瞬时速度。
\end{li}

因为 $h(t)=5t^2$，套用 $x^2$ 的求导结果（常数倍数会原样跟着走，理由下面讲），$h'(t)=10t$。所以 $t=10$ 秒时瞬时速度是 $100$ 米每秒。注意我们一步平均速度都没算，但给出的答案比任何一段平均速度都精确——它说的就是“第 10 秒那一瞬”。

\section{为什么导数有用：局部线性近似}

导数最深刻的含义，不是“一个斜率”，而是“一条曲线在一个点附近的最好直线替身”。

回到 $f(x)=x^2$ 在 $x=1$ 处。切线方程是 $y=2x-1$。我们比较一下曲线和切线在 $x=1.02$ 处的值：曲线给出 $1.02^2=1.0404$，切线给出 $2\times 1.02-1=1.04$。两者相差只有 0.0004！也就是说，在 $x=1$ 附近，如果你懒得算平方，直接用切线 $2x-1$ 去近似 $x^2$，误差小到可以忽略。把镜头怼得足够近，任何光滑曲线都“变成”了直线——第 8 章预告过的“局部像直线”，现在有了精确的度量：那条直线就是切线，它的斜率就是导数。

这件事有一个直接的应用：估算。比如估算 $\sqrt{4.1}$。函数 $g(x)=\sqrt{x}$ 的导数可以算出是 $\frac{1}{2\sqrt{x}}$（极限定义配合分子有理化可得），在 $x=4$ 处导数是 $\frac{1}{4}$。用切线近似：
\[
\sqrt{4.1}\approx \sqrt{4}+\frac{1}{4}\times 0.1=2+0.025=2.025.
\]
真值约是 2.02485，误差不到 0.0002。没有计算器也能得到三位有效数字——导数替我们干的。

值得把这个近似公式单独端详一遍：
\[
f(x_0+\Delta x)\approx f(x_0)+f'(x_0)\,\Delta x.
\]
左边是曲线上真实的点，右边是切线上的点。整件事的结构是：“现在的值”加上“变化率乘以步长”。变化率越大，同样一小步造成的变化就越大——这正和我们日常的直觉吻合：陡坡上的每一步都更费力。但要记住这个近似的边界：它只在 $\Delta x$ 很小时可靠，而且曲线越“弯”（二阶导数越大），近似失效得越快。近似不是糊弄，而是标明了误差的权衡取舍。

\begin{toolbox}
本章新工具：
\begin{itemize}
  \item \textbf{导数}：平均变化率的极限，几何上是切线斜率，物理上是瞬时速度。
  \item \textbf{幂函数求导}：$(x^n)'=nx^{n-1}$（本章算了 $n=2,3$ 两种）。
  \item \textbf{局部线性近似}：$f(x_0+\Delta x)\approx f(x_0)+f'(x_0)\,\Delta x$，离 $x_0$ 越近越准。
  \item \textbf{导数符号的解读}：$f'>0$ 表示函数上升，$f'<0$ 表示下降，$f'=0$ 是可能的“山头、谷底或歇脚处”。
\end{itemize}
\end{toolbox}

\section{函数之间怎么求导：乘法与链式法则}

会算 $x^2$ 和 $x^3$ 还不够，现实中的函数常常是“组装”出来的：两个函数相乘、一个函数套着另一个函数。导数能不能也“组装”？

两个函数相加最简单：$(f+g)'=f'+g'$，因为和的增量就是增量的和，用定义一推就有。常数倍也一样：$(cf)'=cf'$。

但乘法麻烦得多。很多人第一反应是“乘积的导数等于导数的乘积”。错。立刻拿具体例子戳穿：$x^3=x\cdot x^2$，两边我们都认识。$x^3$ 的导数是 $3x^2$；而 $x$ 的导数是 1，$x^2$ 的导数是 $2x$，乘起来只有 $2x$，差得远。正确的规则要复杂一点。

\begin{dingli}[乘积法则]
若 $f$ 与 $g$ 都在 $x$ 处可导，则
\[
(fg)'(x)=f'(x)g(x)+f(x)g'(x).
\]
\end{dingli}

\begin{proof}
考察乘积的差商，并在分子上“借一项、还一项”：
\[
\frac{f(x+h)g(x+h)-f(x)g(x)}{h}
=\frac{f(x+h)g(x+h)-f(x)g(x+h)}{h}
+\frac{f(x)g(x+h)-f(x)g(x)}{h}.
\]
第一项提出 $g(x+h)$，剩 $\frac{f(x+h)-f(x)}{h}\,g(x+h)$；第二项提出 $f(x)$，剩 $f(x)\,\frac{g(x+h)-g(x)}{h}$。令 $h$ 趋于 0：两个差商分别趋于 $f'(x)$ 和 $g'(x)$，而 $g(x+h)$ 趋于 $g(x)$（可导的函数必定连续，这一点可用定义直接验证）。于是极限是 $f'(x)g(x)+f(x)g'(x)$。
\end{proof}

几何上这个公式很形象：把 $f$ 和 $g$ 想象成一个矩形的两条边长，面积就是 $fg$。两条边都变长一点，面积的变化由两块组成——“长边增长”贡献一块，由现在的宽乘以长的变化率；“宽边增长”贡献一块，由现在的长乘以宽的变化率。（严格说还有一个小角，是两块小增量的乘积，在极限里可以忽略。）

验证一下刚才的反例：$x^3=x\cdot x^2$，按乘积法则，导数是 $1\cdot x^2+x\cdot 2x=x^2+2x^2=3x^2$。完全对上了。

再说函数套函数，比如 $y=(x^2+1)^3$。这里面有个“内层”$u=x^2+1$，一个“外层”$u^3$。直觉是：$x$ 动一点，$u$ 跟着动，$u$ 一动，$y$ 再动。变化率应该连乘。

\begin{dingli}[链式法则]
若 $u=g(x)$ 在 $x$ 处可导，$y=f(u)$ 在对应的 $u$ 处可导，则复合函数 $y=f(g(x))$ 在 $x$ 处可导，且
\[
\frac{\mathrm{d}y}{\mathrm{d}x}=\frac{\mathrm{d}y}{\mathrm{d}u}\cdot\frac{\mathrm{d}u}{\mathrm{d}x}.
\]
\end{dingli}

我们不证完整版（需要一点关于“两个极限过程交换”的技术细节），但说明它为什么合理：当 $\Delta x$ 很小时，$\Delta u\approx g'(x)\Delta x$，$\Delta y\approx f'(u)\Delta u$，代进去就是 $\Delta y\approx f'(u)g'(x)\Delta x$。两个近似的误差在取极限后消失。

算一下例子：$y=(x^2+1)^3$。外层 $u^3$ 的导数是 $3u^2$，内层 $x^2+1$ 的导数是 $2x$，链式法则给出
\[
y'=3(x^2+1)^2\cdot 2x=6x(x^2+1)^2.
\]
要知道这有多省事，可以试试把 $(x^2+1)^3$ 展开成 $x^6+3x^4+3x^2+1$ 再逐项求导，结果当然是 $6x^5+12x^3+6x$，因式分解后一模一样——但链式法则一行就写完了，而且还告诉你变化是“哪一环贡献的”。

\section{换一个视角：误差、优化与经济学中的边际}

导数一旦普及，许多原本“试着来”的问题就有了系统解法。这里看四个视角。

\textbf{误差估计。}你量一个正方形的边长，量出 10 厘米，但尺子有误差，真实边长可能是 $10\pm 0.05$ 厘米。面积的误差有多大？面积 $S=x^2$，导数 $S'=2x$，在 $x=10$ 处是 20。用局部线性近似：边长差 0.05，面积约差 $20\times 0.05=1$ 平方厘米。导数告诉你“测量误差会被放大多少倍”——这个倍数就是导数本身。这也是工程师检验一个设计“稳不稳”的第一件事：如果某个量的导数很大，微小的制造误差就会被放大成严重后果。

\textbf{优化。}用 40 米围栏靠墙围一个矩形菜园（墙占一条长边），怎样围面积最大？设宽为 $x$ 米，长就是 $40-2x$ 米，面积
\[
S(x)=x(40-2x)=40x-2x^2.
\]
以前这类题靠配方或画图猜，现在有导数：$S'(x)=40-4x$。导数为正，面积随 $x$ 增加而增加；导数为负，面积开始减少；所以从增到减的转折点一定在导数为零处：$40-4x=0$，$x=10$。此时长为 20 米，最大面积 200 平方米。山顶和谷底的共同特征是“切线水平”，即导数为零——这是找极值的第一张地图（当然还要检查是不是边界或歇脚处，见下面的反例警报）。

\textbf{边际。}经济学家问的不是“一共生产了多少”，而是“再多生产一件，成本增加多少”。成本函数 $C(x)$ 在产量 $x$ 处的导数叫边际成本。一个工厂总成本 $C(x)=1000+20x+0.01x^2$ 元，产量 100 件时的边际成本是 $C'(100)=20+0.02\times 100=22$ 元。这意味着第 101 件产品大约增加 22 元成本——注意是“大约”，因为导数是连续的局部近似，而产品是一件一件离散的。明确适用边界，这是用数学建模的好习惯。

\textbf{运动。}回到本章开头的汽车。路程的导数是速度，那速度的导数是什么？是速度变化的快慢——加速度。电梯启动的一瞬间你感到轻微下坠，转弯时身体被甩向一侧，你感受到的都不是速度本身，而是加速度。设自由落体的下落距离是 $s(t)=4.9t^2$ 米，那么速度是 $v(t)=s'(t)=9.8t$ 米每秒；速度再求一次导数，$v'(t)=9.8$，是一个常数——这正是伽利略在斜面实验里发现的事实：无论多重、无论下落了多久，物体下落的加速度几乎不变。导数连求两次，就把伽利略成箱的实验数据压缩成了一行公式。物理学家把“导数的导数”叫二阶导数，它描述了曲线“弯得有多快”，在第 8 章预告过的曲率概念里还会再见面。

\begin{lianjie}
导数把本书前几章串了起来：第 2 章的斜率是常数函数的特例；第 8 章“曲线局部像直线”的承诺，在本章变成了切线与导数；第 9 章的极限，是导数定义的骨架。往后看，第 12 章会证明“导数为零的点是极值候选点”的严格版本，第 17 章会把“切线”升级成曲面的“切平面”。而在物理、经济、生物里，凡是有“变化率”的地方——人口增长、放射性衰变、传染病扩散——导数都是第一语言。
\end{lianjie}

\begin{shiyan}
\textbf{实验一：望远镜里的切线。}在方格纸上画出 $y=x^2$（取 $x$ 从 $-3$ 到 $3$，每个整点描点连线）。用直尺在 $(2,4)$ 处沿曲线方向轻轻画出一条“只碰不跨”的直线，量它的斜率（数格子算“升多少除以走多少”）。再算一算：理论上应该是 $2x=4$。你的手画结果离 4 有多远？

\textbf{实验二：切线估算。}不许用计算器，利用 $f(x)=x^2$ 在 $x=5$ 处的切线估算 $5.1^2$。切线是 $y=10x-25$，代入 $x=5.1$ 得 26。真实值 $26.01$，你的估算差了 0.01。再用同样方法估 $5.5^2$，看看误差变成多少，体会“离切点越远，近似越差”。

\textbf{实验三：斜率的符号侦探。}随手画一条光滑曲线（比如先上升、再下降、再上升的波浪线），不计算，只在图上的几个点标注该处导数是正、是负还是零。然后同桌交换检查。这个游戏练的是“看见导数”的眼睛。
\end{shiyan}

\section{再看一眼“山顶”：导数为零与极值}

刚才围栏问题里我们用了“极大值处导数为零”。把它写成一条正式的定理。

\begin{dingli}[极值的必要条件]
设 $f$ 在某个区间内部的可导点 $x_0$ 处取得最大值或最小值，则 $f'(x_0)=0$。
\end{dingli}

\begin{proof}
以最大值为例。既然 $f(x_0)$ 是最大，那么对任何非零的 $h$（只要 $x_0+h$ 还在区间里）都有 $f(x_0+h)\le f(x_0)$。当 $h>0$ 时，差商 $\frac{f(x_0+h)-f(x_0)}{h}\le 0$，取极限得 $f'(x_0)\le 0$；当 $h<0$ 时，分母变负，不等式翻转，差商 $\ge 0$，取极限得 $f'(x_0)\ge 0$。一个数既不大于 0 又不小于 0，只能是 0。
\end{proof}

注意这条定理的语言很克制：它只说“在山顶，切线一定水平”，从不说“切线水平的地方一定是山顶”。这个区别至关重要。

\begin{fanli}
有人说：“$f'(x_0)=0$，所以 $x_0$ 处是极大值或极小值。”举 $f(x)=x^3$ 就戳穿了：$f'(x)=3x^2$，在原点处导数为 0，可原点左边函数值为负、右边为正，曲线一路单调上升，原点既不是山顶也不是谷底——它只是我们前面说的“歇脚处”（数学上叫驻点或拐点情形）。

还有第二个坑：定理要求“区间内部”。$f(x)=x$ 在 $[0,1]$ 上的最大值在 $x=1$ 处取得，但那里导数是 1，不是 0。闭区间上的最值还可能藏在端点。所以找最值的完整流程是：算出所有导数为零的点、导数不存在的点，再加上端点，逐个比较函数值——一张候选名单，一个都不能漏。
\end{fanli}

\begin{tuilun}
若可导函数 $f$ 在区间上处处满足 $f'(x)=0$，则 $f$ 是常数函数。
\end{tuilun}

这条推论听起来平凡，却是微积分里最有用的“侦探工具”之一：想知道两个函数是不是只差一个常数，只要算它们的差的导数是不是恒为零。它的严格证明需要第 12 章的中值定理，我们这里先记下它作为悬念。

\section{挑战：亲手算一条“会拐弯的跑道”}

综合本章工具，解一个完整问题。游乐园设计一条滑道，高度（米）随水平距离 $x$（米）变化：$y=x^3-6x^2+9x+1$（$0\le x\le 4$）。问：（1）滑道在哪里最陡地下降？（2）滑道的最高点和最低点在哪里？

先求导。逐项来：$(x^3)'=3x^2$，$(-6x^2)'=-12x$（常数倍法则），$(9x)'=9$，常数项导数为 0。所以
\[
y'=3x^2-12x+9=3(x^2-4x+3)=3(x-1)(x-3).
\]
导数为零的点：$x=1$ 和 $x=3$。看符号：$x<1$ 时两个因子同号（都负），导数为正，滑道上升；$1<x<3$ 时导数为负，下降；$x>3$ 时又为正，上升。所以 $x=1$ 是山顶，$y(1)=1-6+9+1=5$ 米；$x=3$ 是谷底，$y(3)=27-54+27+1=1$ 米。

但题目问的是闭区间 $[0,4]$ 上的最高最低点，端点也要查：$y(0)=1$，$y(4)=64-96+36+1=5$。比较四个值 $\{5,1,1,5\}$：最高 5 米（在 $x=1$ 和 $x=4$），最低 1 米（在 $x=0$ 和 $x=3$）。如果只盯着导数为零的点，就会漏掉端点的答案——这正是上面反例警报的用武之地。

至于“最陡地下降”，那是导数最负的位置。$y'=3(x-1)(x-3)=3(x-2)^2-3$，这是关于 $x=2$ 的开口向上的抛物线，最小值在 $x=2$ 处取得，$y'(2)=-3$。所以滑道在 $x=2$ 米处下坡最猛，斜率 $-3$——恰好是山顶与谷底连线的中点，这个巧合在 $x^3$ 类曲线里普遍成立，原因藏在“导数的导数”里，那是继续攀登的方向。

\begin{pandeng}
本章只是微积分的第一级台阶。继续往上的关键词：\textbf{二阶导数}（导数的变化率，描述凹凸与加速度）、\textbf{中值定理}（连接函数值与导数的桥梁，是“导数为零推出常数”的严格靠山）、\textbf{洛必达法则}（用导数算 $\frac{0}{0}$ 型极限）、\textbf{泰勒展开}（把局部线性近似升级成二次、三次乃至无穷多项的逼近）、\textbf{多元函数与梯度}（曲面上每个点有一整个方向的“斜率集合”）。竞赛里常见的“构造函数求导证不等式”，本质也是本章工具。
\end{pandeng}

\section*{思考题}

\textbf{观察题}
\begin{sikao}
  \item 画一条你自己编的曲线，找出三个导数为零的点。它们都是极大或极小吗？\emph{提示：参考 $x^3$ 在原点附近的样子。}
  \item 函数 $f(x)=|x|$ 在 $x=0$ 处的图像有一个尖角。分别用 $h>0$ 和 $h<0$ 计算差商的极限，你得到什么？\emph{提示：左右两个极限不相等意味着什么？}
\end{sikao}

\textbf{证明题}
\begin{sikao}
  \item 用导数定义证明 $(x^4)'=4x^3$。\emph{提示：先展开 $(x+h)^4$，含 $h^2$ 及以上的项在取极限后全部消失；可以猜想并验证一般规律 $(x^n)'=nx^{n-1}$。}
  \item 用乘积法则证明：常数倍法则 $(cf)'=cf'$ 是乘积法则的特例。\emph{提示：常数函数的导数是多少？}
\end{sikao}

\textbf{挑战题}
\begin{sikao}
  \item 某工厂生产 $x$ 件产品的利润（元）为 $L(x)=-2x^2+80x-300$。求利润最大的产量，并说明为什么“再多生产反而亏钱”。\emph{提示：先求 $L'(x)=0$ 的点，再解释导数符号在该点两侧的变化。}
  \item 证明：在可导的前提下，偶函数（满足 $f(-x)=f(x)$）的导数是奇函数。\emph{提示：对 $f(-x)=f(x)$ 两边用链式法则求导，或回到差商定义分别比较 $f'(-a)$ 与 $f'(a)$。}
\end{sikao}

\section*{通往下一章的门}

这一章我们学会了“从累积量看变化率”：知道路程函数，就能求出每一瞬的速度；知道成本函数，就能求出边际成本。整个过程是一个极限、一个除法。

但把问题反过来问，就没有除法可用了：\emph{如果我只知道汽车每一时刻的速度表读数，能不能算出它一共开了多远？}或者换一个几何版本：\emph{知道曲线 $y=x^2$，能不能算出它和 $x$ 轴之间围成的面积？}直边图形的面积我们会算，可这里的边界是弯的。

直觉暗示：把整段路切成无数极短的小段，每段近似匀速，路程就近似是“速度乘时间”，再把无数小段加起来。段切得越细，近似越准——这又是一个极限，但方向相反：本章是“差商的极限”，下一章是“和的极限”。更令人吃惊的是，这两个方向竟然互为逆运算，像乘法和除法一样互相解开：对累积量求导得到变化率，对变化率“求和”还原出累积量。这个深刻的对偶，就是微积分基本定理，也是下一章的故事。
